\structsection{Заключение}

В выпускной квалификационной работе была исследована задача Smart 3D Packing, представляющая собой прикладную постановку 3D Bin Packing для укладки коробов на паллету. В отличие от абстрактной геометрической задачи, рассматриваемая постановка учитывает ограничения, характерные для складской логистики: габариты паллеты, грузоподъемность, допустимые повороты, опору, хрупкость, штабелируемость и время расчета.

В ходе работы были изучены условие задачи, материалы НИР и курсовой работы, требования к оформлению ВКР, примеры прошлых дипломов и исходный код проекта. По результатам анализа восстановлен pipeline решения: генерация или загрузка JSON-датасета, запуск одного из решателей, валидация жестких ограничений, расчет метрик, сохранение результатов и визуализация укладки.

Были рассмотрены реализованные методы: greedy на основе extreme points и gravity drop, послойный MaxRects-решатель, Large Neighborhood Search, гибрид GAN+GA, branch-and-bound для малых задач и многопаллетное расширение. Для каждого метода описаны основная идея, роль в проекте, преимущества и ограничения.

Отдельное внимание было уделено тому, что все методы сравниваются в единой постановке. Входные данные имеют общий JSON-формат, результаты решателей приводятся к единой структуре ответа, а проверка выполняется одним валидатором. Благодаря этому различия между алгоритмами интерпретируются как различия в качестве построенной укладки, а не как следствие разных правил проверки или разных форматов результатов.

Экспериментальная часть основана на сохраненных результатах репозитория. Сравнение на общих single-pallet наборах показало, что LNS является наиболее эффективным методом по итоговой метрике: 0.7111 на наборе из 100 задач, 0.7434 на наборе из 200 задач и 0.7298 на контрольном наборе. Его преимущество связано с тем, что метод умеет перестраивать уже найденную укладку, исправляя локально неудачные ранние решения greedy.

Greedy показал высокую скорость и идеальный временной балл, но уступил по утилизации объема и покрытию заказа. MaxRects достиг высокой плотности, но потерял качество по хрупкости. GAN+GA подтвердил перспективность гибридного ML-ориентированного поиска, однако в текущей реализации не превзошел LNS. Brute force оказался полезен для малых задач, но плохо масштабируется на более крупные экземпляры.

Проведенный анализ также показал ограничения текущего этапа работы. Эксперименты основаны преимущественно на синтетических данных, физическая модель устойчивости упрощена до проверки площади опоры, а стохастические методы требуют дополнительной оценки разброса результатов по нескольким независимым запускам. Эти ограничения не отменяют полученных выводов, но задают направления, в которых проект может быть усилен.

Практическая ценность полученного решения состоит в том, что проект уже содержит не только отдельные алгоритмы, но и полный контур проверки: генерацию задач, запуск решателей, расчет метрик, хранение результатов и визуализацию. Такая организация позволяет постепенно улучшать отдельные методы, не меняя постановку эксперимента целиком. Например, можно заменить destroy-оператор в LNS, расширить генератор данных или уточнить правило хрупкости, после чего проверить изменение результата тем же benchmark-процессом.

Для дальнейшего развития наиболее важной является проверка на данных, близких к реальным складским заказам. Синтетические сценарии полезны для управляемого сравнения, но промышленная эксплуатация требует учитывать дополнительные ограничения: совместимость товарных групп, требования к разгрузке, температурные зоны и особенности конкретной упаковки. Поэтому следующим этапом может стать сбор обезличенных заказов или создание генератора, параметры которого настраиваются по статистике реального склада.

Содержательно работа показывает, что для 3D Bin Packing в логистической постановке наиболее продуктивен комбинированный подход. Быстрые эвристики нужны для построения стартовых решений и контроля времени, послойные методы дают сильную геометрическую базу, метаэвристики позволяют исправлять локально неудачные решения, а ML-ориентированные методы могут использоваться как источник новых порядков и признаков для поиска. В такой системе каждый подход занимает свою роль, поэтому дальнейшее развитие проекта целесообразно строить не как замену всех методов одним универсальным алгоритмом, а как улучшение общего экспериментального контура.

Таким образом, цель работы достигнута: выполнен анализ предметной области, формализована постановка задачи, изучена реализация проекта, описаны алгоритмы и проведено сравнение методов. Наиболее рациональным основным методом для текущей single-pallet постановки является LNS, а дальнейшее развитие проекта целесообразно связывать с усилением его destroy/repair-операторов, улучшением физической модели устойчивости, расширением генератора данных и развитием многопаллетного режима.
